\section{Discussion}
\label{sec:discussion}

For Markov sequences the structure a denoiser can exploit is not in the score, which diffusion makes non-local, but in
the posterior, which stays a tree. BP turns a learned kernel, a function of two variables, into the exact score at every
noise level, and EM learns the kernel from noisy data with BP as its E-step. With a small neural energy kernel this
needs no choice of family: it ties the best hand-designed family where that family is right and stays below the
networks where it is wrong. Where the physics is known, a kernel with the form of the action does better still and
returns interpretable parameters.

\paragraph{What this suggests for diffusion models.} Learn the local law and infer globally. For data that are Markov
in some ordering, one learned kernel gives the score at every noise level by message passing, and a kernel learned on
short sequences applies unchanged to long ones, at a cost linear in $L$. For data that are not, a kernel can handle the
Markov part exactly and a network the rest: \cref{fig:nonmarkov} puts the scale at which the rest matters at about
$10^{-2}$ nats per site, and the pair state shows that part of it can be absorbed into a larger Markov state. On
hierarchical data, too, BP on a tree is the exact denoiser, and U-Nets can be read as approximating it
\cite{mei2024unets,sclocchi2025phase}.

\paragraph{Limits.} The method needs a known Markov order: on a state of $p$ consecutive sites a grid step costs
$O(\Ng^{p+1})$. A state of dimension $d$ needs a grid of $\Ng^d$ points. At $d=2$ a $41\times41$ grid still works: on a
rotated bivariate Laplace chain the neural kernel beats two-channel networks in 7 of 8 seeds at $N\le512$, though it
beats the linear denoiser of the true covariance only at $N=2048$. Beyond that the messages need another representation
(particles, mixtures, low-rank).
% src: tools/core/read_e67.py (E67, frozen 5 Nov 2026, 6c56fbb): energy2d 0.0532/0.0406/0.0304, bigru2 0.0599/0.0438/0.0344,
% var1 0.0455/0.0381/0.0360, lmmse (true covariance) 0.0351 at N = 128/512/2048; levels t >= 0.3 only.
The index set must be a tree. The grid costs $O(L\Ng^2)$ per sequence, 6 to 20 times a sequence network's forward pass,
and limits how small $t$ can be. The neural kernel's M-step needs a safeguard against stalls (\cref{app:em}). Finally, a
large network closes the gap with enough data: on the Allen--Cahn field the U-Net matches the neural kernel with four
times the data. What we measure is data efficiency, which matters most where scientific data are scarce.
% src: stalls E58 23/576 (4.0%), E55 finished cells 15/208 (7.2%, final lr < 1e-10); E61 with the fix 0/448

\paragraph{Related work.} EM with forward--backward as the E-step is the classical way to fit hidden Markov and
state-space models \cite{baum1970,rabiner1989tutorial,shumway1982em,cappe2005}, and grid representations of
non-Gaussian filtering densities go back to \cite{kitagawa1987}. Here the hidden chain is the clean data, the emission
is the diffusion's own noise, and the fitted model is used for its score. Forward--backward is differentiable
\cite{eisner2016}, which we use only to train with the networks' objective. Higher-order chains are handled as
first-order chains on a larger state or by mixture transition models \cite{raftery1985,berchtold2002,xu2026}. Diffusion
models can be trained from corrupted data \cite{daras2023ambient,daras2024consistent,kawar2024gsure}, also by EM
\cite{rozet2024learning,bai2024em}, but they learn a network score. Analytic and structured denoisers
\cite{kamb2024analytic,kadkhodaie2024generalization,li2024hidden} and BP on hierarchical data
\cite{garnierbrun2024transformers,mei2024unets} ask what structure a denoiser exploits; Markov sequences are the
simplest non-Gaussian case in which the exact denoiser is computable and can be learned from noisy data alone.
